\chapter{Simulation Part Draft}



The latent drift is parameterized as a linear map $\mu(z) = Mz + N$, where $M$ is a learned weight matrix and $N$ a bias vector. For the stochastic differential equation
$$
dz_t = \mu(z_t)\,dt + L(z_t)\,dW_t = (Mz_t + N)\,dt + L(z_t)\,dW_t
$$
to be stable, it is necessary that all eigenvalues of $M$ have strictly negative real parts. This ensures that the latent state $z_t$ is mean-reverting: deviations from the long-run mean are pulled back rather than amplified. If any eigenvalue has a positive real part $\lambda > 0$, the solution grows asymptotically as $e^{\lambda t}$, rendering the process explosive.

In the three-factor specification trained here, the weight matrix $M$ has eigenvalues $\{-1.558,\; 0.642 \pm 2.528i\}$. The complex conjugate pair carries a positive real part of $0.642$, which implies an exponential growth rate of $e^{0.642\,t}$ in the unstable directions. The imaginary component further introduces oscillatory behaviour, explaining the spiralling divergence observed in the simulated paths. Concretely, after a 10-year simulation horizon the latent states reach values on the order of $10^5$, far outside the training distribution.

This instability arises because the current parameterisation places no constraint on the spectrum of $M$ during training. To guarantee mean-reversion, the parameterisation of the drift network must be modified so that the matrix $M$ is constrained to be a stability matrix by construction — for example by writing $M = -V^\top V - \varepsilon I$ for a learned matrix $V$ and small $\varepsilon > 0$, which ensures all eigenvalues are strictly negative. Without such a constraint, the model cannot be reliably used for long-horizon simulation or derivative pricing.

